\section{“Pre-train 没有到头”：scaling law 与数据墙}

当主持人问模型进步速度是否放缓时，姚顺宇明确回答“完全没有”。他的理由不是某个 benchmark 还在涨，而是作为研究员的体感：模型学习东西的能力变强了。以前让模型学会某件事需要很多技巧，现在更重要的是把问题定义清楚、构建合适数据和环境，剩下很多事情会“顺其自然”。

\begin{figure}[H]
\centering
\includegraphics[width=0.88\textwidth]{frame-pretraining.jpg}
\caption{“Pre-train 没有到头”章节：访谈把模型进步放缓与 benchmark 饱和分开讨论。\protect\footnotemark}
\end{figure}
\footnotetext{视频画面时间区间：00:30:10--00:30:38。该帧来自预训练章节中段，用于定位讨论场景。}

\begin{knowledgebox}{读图：预训练章节的视觉证据}
画面仍是双人访谈，说明这里没有官方幻灯片或公式板书。阅读时应把“Pre-train 没有到头”理解为口头判断，需要结合后文的 scaling law、数据墙和用户体验讨论，而不是把单帧画面当作实验证据。
\end{knowledgebox}


\begin{importantbox}{核心判断：benchmark 饱和不等于能力增长停止}
如果指标上限是 100\%，越接近上限，月度增长自然变慢。但用户体验的增长不一定线性对应 benchmark 分数。从 70\% 到 75\% 可能比从 50\% 到 60\% 更有体感，也可能在某些任务上没有体感。真正的问题是：模型是否仍能通过更好的预训练、数据和环境获得新的泛化能力。
\end{importantbox}

\subsection{scaling law 的三种“到头”含义}

姚顺宇把“预训练到头”的可能原因拆成几类。第一，规律本身可能有适用范围，无法无限延展。第二，规律需要的某个条件无法满足，例如高质量数据不够。第三，外部观察者可能只看到某类指标饱和，却误以为底层学习能力饱和。

\noindent\begin{tabularx}{\textwidth}{p{0.25\textwidth}p{0.32\textwidth}X}
\toprule
“到头”说法 & 真实含义 & 应如何验证 \\
\midrule
规律适用范围到头 & scaling law 不再外推 & 需要跨规模实验，而不是单点 benchmark 观察。 \\
数据墙 & 高质量数据无法继续扩展 & 要看数据治理、合成数据、环境数据、交互数据能否补足。 \\
指标饱和 & 某个公开测试接近上限 & 要换更难、更真实、更贴近产品的评估。 \\
体验边际下降 & 用户不再感到显著提升 & 要分析任务分布，而不是只看平均分。 \\
\bottomrule
\end{tabularx}

\begin{knowledgebox}{预训练在这里不只是“继续喂文本”}
访谈中的 pre-training 指向更宽泛的底座学习能力。它包括模型从大规模数据中学到世界知识、代码结构、语言模式、任务先验和抽象能力。后训练、工具使用、强化学习和环境数据会改变模型行为，但底座的学习能力仍决定了很多上限。
\end{knowledgebox}

\begin{warningbox}{“数据墙”不是一句口号}
说数据撞墙，至少要回答三件事：哪类数据撞墙，质量定义是什么，替代数据或环境反馈是否真的不可用。姚顺宇的态度是：未来几个月看不到预训练到头的迹象，但这不等于 scaling law 可以无限外推。
\end{warningbox}

\subsection{本章小结}

预训练是否到头，不能只看公开榜单。姚顺宇的判断强调“模型学习能力”仍在增强，而下一阶段更依赖问题定义、数据构造和环境设计。scaling law 的争论本质上是对可扩展资源、数据质量和评估边界的争论。

